\subsection{\bsq{switch} Statements}\label{sec:SwitchStatements}
This chapter describes how to format a \lstinline|switch|\index{switch} statement together with the related \lstinline|case|\index{case} and \lstinline|default|\index{switch!default} statements.

Beginning with the first preview in Java~12, an alternative syntax for \lstinline|switch|\index{switch} was introduced, so that we have now two (in fact, three) different forms of that construct.

Some sources (e.g. \autocite{ORACLE_DOC_SWITCHEXPRESSIONS}) refer to the different syntax variants by the form of the \lstinline|switch|\index{switch} labels as ``\verb#case L:#'' (\bsq{old}) and ``\verb#case L ->#'' (\bsq{new}). 

\begin{tabular}{ll}
\rule[-1ex]{0pt}{2.5ex} ``\verb#case L:#'' or \bsq{old} \lstinline|switch| & ``\verb#case L ->#'' or \bsq{new} \lstinline|switch| \\ 
\rule[-1ex]{0pt}{2.5ex}
\begin{minipage}{0.48\textwidth}
\begin{lstlisting}[numbers=left]
switch( <selector> )
{
    case <label1>:
        <statements1>; 
        break;
    
    case <label2a>:
    case <label2b>: 
        <statements2>; 
        break;
    
    default: <statements>; break;
}
\end{lstlisting}
\end{minipage} & 

\begin{minipage}{0.48\textwidth}
\begin{lstlisting}
switch( <selector> )
{
    case <label1> -> <statements1>;
    
    
        
    case <label2a>, <label2b> 
        -> <statements2>;


    
    default -> <statments>
}
\end{lstlisting}
\end{minipage} \\ 
\end{tabular} 

%==============================================================================
% End of file!!
%==============================================================================


I will discuss in chapter \tqfullvref{sec:TheSwitchStatement} how and when to use which form.

%------------------------------------------------------------------------------

\subsubsection{The traditional Form of \bsq{switch}}
Originally, a \lstinline|switch|\index{switch}\index{switch!statement} statement in \Java looked like this (the ``\verb#case L:#'' variant):
\begin{lstlisting}[numbers=left]
switch( <selector> )
{
    case <switchlabel1>:
        <statements1>;
        // Falls through!

    case <switchlabel2>:
        <statements2>;
        break;

    case <switchlabel3>: <singleStatement>; break;

    case <switchlabel4>:
    {
        <statements>;
        break;
    }

    case <switchlabel5>: // Also a fall-through
    case <switchlabel6>:
    {
        <statements>;
        break;
    }

    default:
       <statements>;
       break;
}
\end{lstlisting}

First, these coding conventions make the blank line above each \lstinline|case|\index{case} (except the very first one) and above \lstinline|default| mandatory!
  
If a \lstinline|case| falls through, like in lines~3 to 9, a comment \bdq{\lstinline|// Falls through!|} (or something similar\footnote{Linting tools as well as most IDEs issue a warning in case of a fall-through. Usually, these could be switched off by adding a comment or an annotation to the location of the fall-through. This would be a sufficient replacement for the comment.}) as in line~5 of the sample is required! But basically, such constructs should be avoided when possible, because it also forces that the sequence of the \lstinline|case| clauses cannot be changed, and that fact requires an additional comment to the \lstinline|switch| itself. In fact here you should consider to ignore the \bdq{DRY~Principle}\index{DRY Principle} – we will see this later – and repeat the code for the second \lstinline|case| on the first one:
\begin{lstlisting}
// BETTER
switch( <selector> )
{
    case <switchlabel1>:
        <statements1>;
        <statements2>;
        break;

    case <switchlabel2>:
        <statements2>;
        break;

    default: …
}
\end{lstlisting}
Or you put the code that is common for both cases into a \lstinline|private| method and call that from both branches. 

Different to the case shown in lines~3 to 9, no comment is required in the case shown in lines~19 and 20: here we have two \lstinline|case| clauses that triggers \textit{exactly the same} action, and we can even omit the blank line between them (and also the comment from line~19).

For longer \lstinline|switch| statements it is highly recommended to use a label with the \lstinline|break| statement\footnote{In fact, I recommend to always use a label with \lstinline|break|; see chapter \tqfullref{sec:LabelsAndBreakStatements} on this topic.}; this can look like this:
\begin{lstlisting}
final Color color;
ColorSelector: switch( colorIndex )
{
    case 1: color = RED; 
        break ColorSelector;
        
    case 2: color = BLUE; 
        break ColorSelector;
        
    case 3: color = YELLOW; 
        break ColorSelector;
        
    case 4: color = GREEN; 
        break ColorSelector;
        
    default: color = NONE; 
        break ColorSelector;
}   //  ColorSelector:
\end{lstlisting}

The \lstinline|break| in the \lstinline|default| branch is redundant, but it prevents a fall-through error if later another \lstinline|case| branch is added – although it should be assured that the \lstinline|default| branch is always the last branch in the \lstinline|switch| statement. Of course no \lstinline|break| is required when the branch is left by throwing an exception.\footnote{A \lstinline|break| would be also obsolete when the branch is left through a \lstinline|return| statement, but in that case, the \lstinline|return| would not be the last statement in the method, and this is another requirement from this coding conventions!}

%==============================================================================
% End of file!!
%==============================================================================

\subsubsection{The new Form of \bsq{switch}}\label{sec:NewSwitch}
The alternative syntax for \lstinline|switch|\index{switch} comes in two flavours (as \textit{statement}\index{switch!statement} and \textit{expression}\index{switch!expression}).

The new \lstinline|switch|\index{switch!statement} statement looks like this:

\begin{lstlisting}[numbers=left]
// switch statement
switch( <selector> )
{
    case <switchlabel1> -> <singleStatement>;

    case <switchlabel2>, <switchlabel3> -> <singleStatement>;

    case <switchlabel4> ->
    {
        <statements>;
    }

    case <switchlabel4>, <switchlabel5> ->
    {
        <statements>;
    }

    default -> <singleStatement>;
}
\end{lstlisting}

And this is the new \lstinline|switch|\index{switch!expression} expression:

\begin{lstlisting}[numbers=left]
// switch expression
var result = switch( <selector> )
{
    case <switchlabel1> -> <expression>;

    case <switchlabel2>, <switchlabel3> -> <expression>;

    case <switchlabel4> ->
    {
        <statements>;
        yield <value>;
    }

    case <switchlabel4>, <switchlabel5> ->
    {
        <statements>;
        yield <value>;
    }

    default -> <expression>;
}
\end{lstlisting}

Of course the \lstinline|default| in line~18 can be followed by a statement block as the \lstinline|case| in line~8, and the \lstinline|default| in line~40 can have a complex expression like the \lstinline|case| in line~28.

As this format does not allow a fall-through in cases where the branches for two or more labels are doing the same stuff, it is possible here to have more than one switch label per \lstinline|case|, as shown in lines~6, 13, 26 and 34.

If used with pattern matching (refer to \autocite{ORACLE_DOC_PATTERNMATCHING}), a \lstinline|switch| looks like this:
\begin{lstlisting}[numbers=left]
var selector = <anObject>
// switch statement
switch( selector )
{
    case null -> <singleStatement>;
    
    case <class1> <name1> -> <singleStatement>;

    case <class2> <name2> ->
    {
        <statements>;
    }

    default -> <singleStatement>;
}

// switch statement
var result = switch( selector )
{
    case null -> <expression>;

    case <class1> <name1> -> <expression>;

    case <class2> <name2> ->
    {
        <expressions>;
        yield <value>;
    }

    default -> <expression>;
}
\end{lstlisting}

Same as for the \lstinline|default| branch, the \lstinline|case null| branch can be a statement block or a complex expression.

And an expression can also be always a \lstinline|throw|.


%==============================================================================
% End of file!!
%==============================================================================

\subsubsection{\bsq{switch} and Pattern Matching}\label{sec:PatternMatchingSwitch}
Pattern matching\index{pattern matching} (refer to \autocite{ORACLE_DOC_PATTERNMATCHING}\index{pattern matching}) is covered in more detail in chapter \tqfullvref{sec:PatternMatching}. Here I want to discuss only the basic formatting of the \lstinline|switch|\index{switch} when used with pattern matching.

\keeplisting{14}
A \lstinline|switch| statement\index{switch!statement} for the ``\verb#case L:#'' variant would look like this:
\index{switch!case}\index{switch!default}\index{switch!break}
\begin{lstlisting}
var selector = <anObject>
switch( selector )
{
    case <class1> <name1>: <singleStatement>;

    case <class2> <name2>:
    {
        <statements>;
    }

    default: <singleStatement>;
}
\end{lstlisting}

\keeplisting{15}
For the ``\verb#case L:#'' variant, a \lstinline|switch| statement\index{switch!statement} is written as below:
\index{switch!case}\index{switch!default}
\begin{lstlisting}
var selector = <anObject>
switch( selector )
{
    case null -> <singleStatement>;
    
    case <class1> <name1> -> <singleStatement>;

    case <class2> <name2> ->
    {
        <statements>;
    }

    default -> <singleStatement>;
}
\end{lstlisting}

\keeplisting{16}
And an ``\verb#case L->#'' variant \lstinline|switch| expression\index{switch!expression} would have this form:
\index{switch!case}\index{switch!default}\index{switch!yield}
\begin{lstlisting}
var selector = <anObject>
var result = switch( selector )
{
    case null -> <expression>;

    case <class1> <name1> -> <expression>;

    case <class2> <name2> ->
    {
        <expressions>;
        yield <value>;
    }

    default -> <expression>;
}
\end{lstlisting}

Same as for the \lstinline|default|\index{switch!default} branch, the \lstinline|case null| \index{switch!case null} branch can be a statement block or a complex expression.

And an expression can also be always a \lstinline|throw|\index{throw} statement.



%==============================================================================
% End of file!!
%==============================================================================

\subsubsection{General formatting Rules for \bsq{switch}}

The general use of blanks, round brackets, curly braces and linefeeds was already discussed in \tqfullvref{sec:WhiteSpace}. The curly braces after the initial \lstinline|switch|\index{switch} are required by the Java syntax and not optional. In case of a \lstinline|switch| statement\index{switch!statement} the opening and closing curly braces are on the same indentation level than the \lstinline|switch| keyword, while for an expression\index{switch!expression} they are aligned with the begin of the respective statement. In both cases, the curly braces are on lines of their own.

The branches (\lstinline|case|\index{switch!case} or \lstinline|default|\index{switch!default}) are indented one level and will be separated from each other by a blank line, to improve the readability.

Although not required (at least not for the ``\verb#case L:#'' variant), curly braces should be used for complex branch statements/expressions, also for readbility. The curly braces for the branches are indented to the same level to the same level as the \lstinline|case|/\lstinline|default|\index{switch!case}\index{switch!default} keywords, while the branch code itself is indented one additional level.

The code for a branch should be short; I suggest that a branch should be limited to less than ten lines. If it is just onle line and it is short enough, the whole branch should be written as a single line, in particular for the ``\verb#case L ->#'' variant:

\index{switch!case}\index{switch!break}
\begin{lstlisting}
case <switchlabel>: <singleStatement>; break;
\end{lstlisting}

\index{switch!case}
\begin{lstlisting}
case <switchlabe> -> <singleStatement>;
\end{lstlisting}

\index{switch!case}
\begin{lstlisting}
case <switchLabel> -> <expression>;
\end{lstlisting}

If necessary, create a private method with the branch code and call that from the branch.

The branches should be ordered somehow logically (alphabetically, by ordinal number,~…), and an existing \lstinline|case null|\index{switch!case null} branch should be the first, while the \lstinline|default|\index{switch!default} branch should be the last one. This makes it easier to apply changes to the \lstinline|switch|\index{switch} later.

\subsubsection{Principles}
\begin{enumerate}[label=P\arabic*.]
\item{No blanks between the \lstinline|switch|\index{switch} and the opening round bracket. Refer to \tqfullvref{sec:ParameterParenthesis} on how to place the white space.}
\item{Curly braces are recommended when the branch statement has more than one line.}
\item{The curly braces for a branch statement are on the same indentation level as the keyword itself.

The branch statements itself are indented one level relative to the curly braces.}
\item{The code for a branch should be short.}
\item{Apply a logical order to the branches.

Place \lstinline|case null|\index{switch!case null} as first, and \lstinline|default|\index{switch!default} as last branches.}
\item{The blank line separating the branches is mandatory.}
\item{No empty branch! A branch must contain at least a comment.

An exception for the ``\verb#case L:#'' variant are branches with the same code.}
\item{Avoid ‘fall-through’!}
\item{Although it looks like one, a \lstinline|switch| expression\index{switch!expression} is not a lambda\index{lambda}.

That means that local variables that are referenced in a \lstinline|switch| expression do not have to be effectively final.}
\end{enumerate}

%==============================================================================
% End of file!!
%==============================================================================


%==============================================================================
% End of file!!
%==============================================================================
